% Stand-alone pedagogical notes (LaTeX)
% Complement to: "Recursive Cognition and Revealed Preferences: A Structural Estimator for Synthetic Cultural Agents" (Feb 9, 2026)

\documentclass[11pt]{article}

\usepackage[margin=1in]{geometry}
\usepackage[T1]{fontenc}
\usepackage{lmodern}
\usepackage{microtype}
\usepackage{amsmath, amssymb, amsthm, mathtools}
\usepackage{booktabs, tabularx, array}
\usepackage{enumitem}
\usepackage{graphicx}
\usepackage{xcolor}
\usepackage[hidelinks]{hyperref}
\usepackage{tcolorbox}
\usepackage{tikz}

% --- Styling helpers ---
\definecolor{ink}{RGB}{20,20,20}
\definecolor{muted}{RGB}{90,90,90}
\definecolor{accent}{RGB}{30,90,160}
\definecolor{light}{RGB}{245,247,250}

\setlist[itemize]{leftmargin=*, topsep=3pt, itemsep=2pt}
\setlist[enumerate]{leftmargin=*, topsep=3pt, itemsep=2pt}

\newcolumntype{Y}{>{\raggedright\arraybackslash}X}
\newcommand{\R}{\mathbb{R}}
\newcommand{\E}{\mathbb{E}}
\newcommand{\KL}{\mathrm{KL}}
\newcommand{\1}{\mathbf{1}}
\newcommand{\sigmoid}{\sigma}
\newcommand{\dd}{\,\mathrm{d}}

\theoremstyle{plain}
\newtheorem{proposition}{Proposition}
\newtheorem{theorem}{Theorem}
\newtheorem{corollary}[theorem]{Corollary}
\newtheorem{lemma}{Lemma}
\theoremstyle{definition}
\newtheorem{definition}{Definition}
\theoremstyle{remark}
\newtheorem{remark}{Remark}

\newtcolorbox{keypoint}{colback=blue!5!white,colframe=blue!75!black,title=Key Point}
\newtcolorbox{econintuition}{colback=green!5!white,colframe=green!75!black,title=Economic Intuition}
\newtcolorbox{firstprinciples}{colback=orange!5!white,colframe=orange!75!black,title=First Principles}


\title{\vspace{-1.0em}\textbf{DPO as Bradley--Terry MLE: A First-Principles Bridge Between\\
Stochastic Control and Structural Estimation}}
\author{Pedagogical companion notes for first-year PhD economics}
\date{(Complement to the attached report; self-contained lecture notes)}

\begin{document}
\maketitle
\vspace{-1.0em}

\begin{tcolorbox}[colback=light, colframe=accent, title=What this document is for]
These notes are designed to complement the attached report by building the equivalence
\[
\text{DPO policy estimation } \Longleftrightarrow \text{Bradley--Terry MLE}
\]
from first principles, with special attention to \emph{what objects are being optimized or estimated},
what assumptions are required for the equivalence to hold, and how the same algebra reappears in
recursive methods (soft Bellman equations) and structural dynamic discrete choice (Hotz--Miller inversions).
\end{tcolorbox}

\paragraph{Main takeaway (in one sentence).}
If (i) preferences are generated by a Bradley--Terry random utility model in \emph{reward differences} and
(ii) the policy that generated those rewards is the optimizer of a \emph{KL-regularized} control problem
around a reference policy, then the \emph{likelihood} of observed preference labels can be written purely
in terms of \emph{policy log-probability ratios}, making DPO exactly the maximum likelihood estimator
for a well-defined paired-comparison model.

\tableofcontents

\section{Objects and ontological hygiene: policy vs.\ estimator}
A lot of confusion in this literature comes from silently switching between \emph{mathematical objects}:
functions, correspondences, probability measures, and data-to-parameter maps. Let's pin them down.

\subsection{Policies are probability kernels (not parameters)}
Fix a measurable context space $\mathcal{X}$ (prompts, states, environments). For each $x\in\mathcal{X}$,
let $\mathcal{Y}(x)$ be the set of feasible actions/outcomes (candidate completions, action sequences, etc.).
A \textbf{policy} is a conditional distribution over $\mathcal{Y}(x)$:
\[
\pi(\cdot\mid x)\in \Delta(\mathcal{Y}(x)).
\]
Formally, $\pi$ is a Markov kernel: for each $x$, a probability measure on $\mathcal{Y}(x)$, and for each measurable $A\subseteq \mathcal{Y}(x)$, the map $x\mapsto \pi(A\mid x)$ is measurable.

\subsection{Estimators map data to parameters (and thereby to policies)}
An \textbf{econometric estimator} is a map from a dataset $D$ to an element of a parameter space $\Theta$,
\[
\hat\theta:\ \mathcal{D}\to\Theta.
\]
If a policy is parameterized by $\theta$---written $\pi_\theta$---then composing gives an \emph{estimated policy}
$\hat\pi(\cdot\mid x)\coloneqq \pi_{\hat\theta(D)}(\cdot\mid x)$.

\begin{remark}[Answering the first-principles question]
An \emph{optimal policy} and an \emph{estimator} are not the same kind of object.
A policy is a conditional distribution; an estimator is a data-dependent rule returning an element of a parameter
space (or a function/correspondence). They become intimately related when we (a) posit a parametric family
$\{\pi_\theta\}$ and (b) choose $\hat\theta$ by solving an optimization problem. Then the \emph{output} of
the estimator is a policy---but the estimator itself is still a data-to-parameter map.
\end{remark}

\begin{definition}[Policy Function]
A \textbf{policy function} $\pi: \mathcal{S} \to \Delta(\mathcal{A})$ maps states to probability distributions over actions. It is a \emph{decision rule} that characterizes optimal behavior given primitives $(u, \beta, \delta)$.
\end{definition}

\begin{definition}[Structural Estimator]
A \textbf{structural estimator} $\hat{\theta}$ is a statistic computed from data $\mathcal{D}$ that recovers parameters of a behavioral model (e.g., preferences, technology).
\end{definition}

\noindent \textbf{When are they the same?}

By the principle of \emph{revealed preference}, observed choices \emph{reveal} the underlying policy. If we observe $(s_i, a_i)$ pairs and posit a parametric policy $\pi_\theta(a|s)$, then:

$$
\hat{\theta}_{\text{MLE}} = \arg\max_\theta \sum_{i=1}^N \log \pi_\theta(a_i | s_i)
$$

is both:
\begin{itemize}
\item An \textbf{estimator} of the policy parameters $\theta$
\item A characterization of the \textbf{optimal policy} under the assumption that observed behavior is optimal
\end{itemize}

\begin{econintuition}
\textbf{The Bridge:} Structural econometrics \emph{assumes} optimality and uses data to recover the policy. Dynamic programming \emph{derives} the policy from primitives. DPO shows these are inverses of each other: the policy that explains preference data under BT is structurally equivalent to the policy that maximizes KL-regularized expected reward.
\end{econintuition}


\subsection{Two ``problems'' whose solutions we'll meet}
We will study two models that live on the same $(x,y)$ space but tell different stories:
\begin{itemize}
  \item A \textbf{positive (statistical)} model: how preference labels are generated from latent utilities
  (Bradley--Terry / random utility).
  \item A \textbf{normative (control)} model: which policy would be chosen if an agent maximized reward
  subject to an information/control cost (KL-regularized planner).
\end{itemize}
The equivalence arises because the KL-regularized optimizer has a \emph{log-ratio inversion} that turns reward
differences into policy log-probability ratio differences. That makes the statistical likelihood depend only on the policy.

\subsection{What is ``Structural'' About DPO?}

A method is \textbf{structural} if it:
\begin{enumerate}
\item Posits economic primitives (preferences, technology, information)
\item Derives testable implications (choice probabilities, moments)
\item Recovers primitives from data via identification arguments
\end{enumerate}

\noindent \textbf{DPO is structural because:}

\begin{itemize}
\item \textbf{Primitives:} Latent reward function $r^*(x,y)$, reference policy $\pi_{\text{ref}}$, rationality parameter $\beta$
\item \textbf{Implications:} Preference probabilities $P(y_w \succ y_\ell | x) = \sigma(\beta^{-1}(r^*(x,y_w) - r^*(x,y_\ell)))$
\item \textbf{Identification:} Log-ratio inversion recovers reward \emph{differences} from observed preferences
\end{itemize}

Contrast with \textbf{reduced-form} methods like supervised fine-tuning, which directly fit $\pi_\theta(y|x)$ to observed $(x,y)$ pairs without a behavioral model.

\newpage
\section{The two model cards (fixed structure, side-by-side comparability)}
To prevent ``symbol drift,'' we use a common model-card template. Each card answers: purpose, primitives,
data, assumptions, estimand, identification, computation, and failure modes.

\vspace{0.5em}
\begin{tcolorbox}[colback=light, colframe=accent, title=Common notation]
\begin{itemize}
\item $x\in\mathcal{X}$: context (prompt / state).
\item $y\in\mathcal{Y}(x)$: candidate completion / action / outcome.
\item $D=\{(x_i, y_i^{(w)}, y_i^{(\ell)})\}_{i=1}^N$: preference comparisons where $y^{(w)}$ is chosen over $y^{(\ell)}$ given $x$.
\item $\pi_{\mathrm{ref}}(\cdot\mid x)$: reference policy (baseline distribution).
\item $\beta>0$: scale/temperature parameter.
\end{itemize}
\end{tcolorbox}

\subsection{Model card A: Bradley--Terry random utility model (paired comparisons)}
\begin{tcolorbox}[colback=white, colframe=accent, title=Model card A: Bradley--Terry (BT) / logistic paired comparison]
\textbf{Purpose.} A tractable probabilistic model for noisy pairwise preference labels.

\medskip
\textbf{Primitives.}
\begin{itemize}
\item A real-valued score (latent utility / reward) $r^\star:\mathcal{X}\times\mathcal{Y}\to\R$.
\item Logistic link $\sigmoid(u)\coloneqq(1+e^{-u})^{-1}$.
\end{itemize}

\medskip
\textbf{Data.} Observations $\{(x_i, y_i^{(w)}, y_i^{(\ell)})\}_{i=1}^N$ where $y_i^{(w)}$ is labeled preferred.

\medskip
\textbf{Model (likelihood).}
\[
\Pr\big(y^{(w)} \succ y^{(\ell)} \mid x\big)=\sigmoid\!\Big(r^\star(x,y^{(w)})-r^\star(x,y^{(\ell)})\Big).
\]

\medskip
\textbf{Estimand.} Utility \emph{differences} (or parameters indexing them). Full levels of $r^\star(x,\cdot)$ are not identified.

\medskip
\textbf{Identification / invariances.}
\begin{itemize}
\item (Additive) For each fixed $x$, $r^\star(x,\cdot)$ is identified only up to $r^\star(x,y)\mapsto r^\star(x,y)+c(x)$.
\item (Scale) If $\beta$ is introduced as a separate preference-noise scale, only the ratio $r^\star/\beta$ matters absent normalization.
\end{itemize}

\medskip
\textbf{Computation.} Convex (in linear scores) logistic regression on differences; nonconvex if $r^\star$ is
neural/LLM-parameterized.

\medskip
\textbf{Failure modes.} Context effects, non-transitivity, rater heterogeneity, dependence on the \emph{presented}
choice set, or systematic violations of the logit error structure imply misspecification. The MLE remains an
$M$-estimator for the best-fitting logit model, but ``utility'' becomes an approximation.
\end{tcolorbox}

\subsection*{Microfoundation: Random Utility Maximization}

The Bradley-Terry model has a clean microfoundation in random utility theory:

\begin{theorem}[Luce-McFadden]
Assume agent $i$ has latent utility:
$$
U_i(x,y) = r^*(x,y) + \varepsilon_y
$$
where $\{\varepsilon_y\}_{y \in \mathcal{Y}(x)}$ are i.i.d.\ Type-I Extreme Value (Gumbel) random variables with CDF $F(\varepsilon) = \exp(-\exp(-\varepsilon))$.

Then:
$$
P(y_1 \succ y_2 | x) = P(U(x,y_1) \geq U(x,y_2)) = \sigma(r^*(x,y_1) - r^*(x,y_2))
$$
\end{theorem}

\begin{proof}[Proof Sketch]
\begin{align*}
P(y_1 \succ y_2) &= P(r^*(x,y_1) + \varepsilon_1 \geq r^*(x,y_2) + \varepsilon_2) \\
&= P(\varepsilon_2 - \varepsilon_1 \leq r^*(x,y_1) - r^*(x,y_2))
\end{align*}

The difference of two i.i.d.\ Gumbel random variables has a logistic distribution:
$$
P(\varepsilon_2 - \varepsilon_1 \leq z) = \frac{e^z}{1 + e^z} = \sigma(z)
$$
\end{proof}

\subsection*{Maximum Likelihood Estimation}

\textbf{Data:} $\mathcal{D} = \{(x_i, y_i^{(w)}, y_i^{(\ell)})\}_{i=1}^N$ where $y^{(w)}$ was preferred to $y^{(\ell)}$ given context $x$.

\noindent \textbf{Parametric Model:} Posit $r^*(x,y) = r_\theta(x,y)$ for some parameter vector $\theta$.

\noindent \textbf{Log-Likelihood:}
$$
\ell(\theta; \mathcal{D}) = \sum_{i=1}^N \log \sigma(r_\theta(x_i, y_i^{(w)}) - r_\theta(x_i, y_i^{(\ell)}))
$$

\noindent \textbf{MLE:}
$$
\hat{\theta}_{\text{MLE}} = \arg\max_\theta \ell(\theta; \mathcal{D})
$$

\subsection*{Identification}

\begin{proposition}[Identification Up to Location]
For each context $x$, the reward function $r^*(x, \cdot)$ is identified only up to an additive constant $c(x)$, since preference probabilities depend only on differences:
$$
\sigma(r^*(x,y_1) - r^*(x,y_2)) = \sigma((r^*(x,y_1) + c(x)) - (r^*(x,y_2) + c(x)))
$$
\end{proposition}

\textbf{Implication:} We can never recover absolute rewards from preference data alone—only reward \emph{rankings} within each context.

\subsection*{Economic Interpretation}

\begin{itemize}
\item \textbf{$r^*(x,y)$}: Systematic utility component
\item \textbf{$\varepsilon_y$}: Unobserved heterogeneity, measurement error, or ``trembles''
\item \textbf{Gumbel assumption}: IIA (Independence of Irrelevant Alternatives)—adding a third option doesn't change relative probabilities of first two
\end{itemize}

\subsection{Model card B: KL-regularized planner and the tilted Gibbs (softmax) policy}
\begin{tcolorbox}[colback=white, colframe=accent, title=Model card B: KL-regularized (entropy) control around a reference policy]
\textbf{Purpose.} A normative model of policy choice trading off expected reward and deviation from a baseline
distribution (information cost / control cost / maximum entropy regularization).

\medskip
\textbf{Primitives.}
\begin{itemize}
\item A reward (or utility) function $r^\star(x,y)$.
\item A reference policy $\pi_{\mathrm{ref}}(\cdot\mid x)$ with $\pi_{\mathrm{ref}}(y\mid x)>0$ on the relevant support.
\item A regularization parameter $\beta>0$ governing the reward--KL tradeoff.
\end{itemize}

\medskip
\textbf{Problem (pointwise in $x$).} Choose $\pi(\cdot\mid x)$ to solve
\begin{equation}\label{eq:klproblem}
\max_{\pi(\cdot\mid x)}\ \E_{y\sim \pi(\cdot\mid x)}[r^\star(x,y)]\;-\;\beta\,\KL\!\big(\pi(\cdot\mid x)\,\|\,\pi_{\mathrm{ref}}(\cdot\mid x)\big)
\quad\text{s.t. } \sum_{y}\pi(y\mid x)=1.
\end{equation}
(Replace sums by integrals if $\mathcal{Y}(x)$ is continuous; the implemented estimand is typically
finite-candidate.)

\medskip
\textbf{Solution (tilted Gibbs / softmax).}
If $r^\star(x,\cdot)$ is bounded (or integrable) and $\pi_{\mathrm{ref}}$ has full support, the objective is strictly
concave and the unique maximizer satisfies
\begin{equation}\label{eq:gibbs}
\pi^\star(y\mid x)=\frac{\pi_{\mathrm{ref}}(y\mid x)\exp(r^\star(x,y)/\beta)}{Z(x)},\qquad
Z(x)=\sum_{y'}\pi_{\mathrm{ref}}(y'\mid x)\exp(r^\star(x,y')/\beta).
\end{equation}

\medskip
\textbf{Key inversion (log-ratio identity).}
Taking logs of \eqref{eq:gibbs} yields, for any $y_1,y_2$,
\begin{equation}\label{eq:logratio}
r^\star(x,y_1)-r^\star(x,y_2)=\beta\Big[\log\frac{\pi^\star(y_1\mid x)}{\pi_{\mathrm{ref}}(y_1\mid x)}-\log\frac{\pi^\star(y_2\mid x)}{\pi_{\mathrm{ref}}(y_2\mid x)}\Big].
\end{equation}
The unknown normalizer $\beta\log Z(x)$ cancels in differences.

\medskip
\textbf{Interpretation.} $\pi_{\mathrm{ref}}$ acts like a baseline ``availability'' (or prior). The KL term penalizes information/control
used to steer the policy away from the baseline. $\beta$ plays the role of a shock scale / rational inattention parameter.

\medskip
\textbf{Failure modes.} If $\pi_{\mathrm{ref}}(y\mid x)=0$ on outcomes that matter, log-ratios diverge and the model breaks.
If the control cost is not KL (e.g., other divergences), the closed-form \eqref{eq:gibbs} and inversion \eqref{eq:logratio} change.
\end{tcolorbox}

\subsection*{The Control Problem}

\begin{definition}[RLHF-KL Objective]
Given a reference policy $\pi_{\text{ref}}: \mathcal{X} \to \Delta(\mathcal{Y})$ and a reward function $r^*: \mathcal{X} \times \mathcal{Y} \to \mathbb{R}$, the KL-regularized optimization problem is:

$$
\max_{\pi} \mathbb{E}_{x \sim \mathcal{D}, y \sim \pi(\cdot|x)} [r^*(x,y)] - \beta \cdot D_{\text{KL}}(\pi(\cdot|x) \| \pi_{\text{ref}}(\cdot|x))
$$

where the KL divergence is:
$$
D_{\text{KL}}(\pi \| \pi_{\text{ref}}) = \sum_{y \in \mathcal{Y}(x)} \pi(y|x) \log \frac{\pi(y|x)}{\pi_{\text{ref}}(y|x)}
$$
\end{definition}

\subsection*{Interpretation}

\noindent \textbf{Why KL regularization?}

\begin{enumerate}
\item \textbf{Robustness:} Prevent overfitting to potentially noisy human preferences
\item \textbf{Fluency:} Stay close to the pre-trained language model distribution
\item \textbf{Computational:} Avoid extreme probability ratios that cause training instability
\item \textbf{Rational Inattention:} $\beta^{-1}$ parameterizes the cost of deviating from the reference
\end{enumerate}

\begin{econintuition}
Think of $\pi_{\text{ref}}$ as a ``default'' or ``baseline'' behavioral pattern (e.g., habitual consumption). The agent optimizes rewards but faces a psychological or computational cost (measured by KL divergence) of deviating from this default. As $\beta \to 0$, the agent becomes infinitely rational (ignores the reference). As $\beta \to \infty$, the agent is boundedly rational (sticks to the reference).
\end{econintuition}

\subsection*{Analytical Solution: The Tilted Gibbs Distribution}

\begin{theorem}[Closed-Form Optimal Policy]
The unique optimal policy for the RLHF-KL problem is:

$$
\pi^*(y|x) = \frac{1}{Z(x)} \pi_{\text{ref}}(y|x) \exp\left(\frac{r^*(x,y)}{\beta}\right)
$$

where $Z(x) = \sum_{y'} \pi_{\text{ref}}(y'|x) \exp(r^*(x,y')/\beta)$ is the partition function.
\end{theorem}

\begin{proof}
Form the Lagrangian:
\begin{align*}
\mathcal{L}(\pi, \lambda) &= \sum_y \pi(y|x) r^*(x,y) - \beta \sum_y \pi(y|x) \log \frac{\pi(y|x)}{\pi_{\text{ref}}(y|x)} + \lambda \left(\sum_y \pi(y|x) - 1\right)
\end{align*}

First-order condition w.r.t.\ $\pi(y|x)$:
$$
0 = r^*(x,y) - \beta \left(\log \frac{\pi(y|x)}{\pi_{\text{ref}}(y|x)} + 1\right) + \lambda
$$

Rearrange:
$$
\log \pi(y|x) = \frac{r^*(x,y)}{\beta} + \log \pi_{\text{ref}}(y|x) + \frac{\lambda - \beta}{\beta}
$$

Exponentiate and normalize using $\sum_y \pi(y|x) = 1$ to get the stated form.
\end{proof}

\subsection*{The Log-Ratio Inversion (Key Step)}

\begin{corollary}[Reward Differences as Log-Ratios]
Taking logs of the optimal policy:
$$
\log \pi^*(y|x) = \frac{r^*(x,y)}{\beta} + \log \pi_{\text{ref}}(y|x) - \log Z(x)
$$

Therefore, for any two alternatives $y_1, y_2$:
\begin{align*}
r^*(x,y_1) - r^*(x,y_2) &= \beta \left[\log \pi^*(y_1|x) - \log \pi_{\text{ref}}(y_1|x)\right] \\
&\quad - \beta \left[\log \pi^*(y_2|x) - \log \pi_{\text{ref}}(y_2|x)\right]
\end{align*}

\textbf{The partition function cancels in differences!}
\end{corollary}

\begin{keypoint}
This is the \textbf{inversion formula}: reward differences can be expressed as log-ratios of the optimal policy relative to the reference, without needing to know $Z(x)$.
\end{keypoint}

\newpage
\section{The equivalence: assumptions, necessity/sufficiency, and the algebra}
We now isolate the exact conditions under which ``DPO is Bradley--Terry MLE'' is literally true.

\subsection{Assumptions (minimal and explicit)}
\begin{enumerate}[label=(A\arabic*)]
\item \textbf{BT data-generating process.} Conditional on $(x,y_1,y_2)$, the probability that a labeler prefers $y_1$ to $y_2$ is
$\sigmoid(r^\star(x,y_1)-r^\star(x,y_2))$ for some latent $r^\star$.
\item \textbf{KL-optimality relation.} The policy that ``implements'' those rewards is the optimizer of the KL-regularized planner
\eqref{eq:klproblem} with reference $\pi_{\mathrm{ref}}$ and temperature $\beta$, giving \eqref{eq:gibbs}.
\item \textbf{Support compatibility.} $\pi_{\mathrm{ref}}(y\mid x)>0$ wherever the policy family (and the data) place mass, so that
log-ratios are well-defined.
\item \textbf{Estimation family.} We fit a parametric policy $\pi_\theta(\cdot\mid x)$ by maximum likelihood under the implied model.
(Correct specification $\exists\theta^\star:\pi_{\theta^\star}=\pi^\star$ is \emph{not} needed for objective equivalence, only for
``recover the true policy'' statements.)
\end{enumerate}

\subsection{Proposition: DPO objective equals BT log-likelihood after substitution}
\begin{proposition}[DPO $\equiv$ BT MLE for log-ratio scores]\label{prop:main}
Under (A1)--(A4), the BT likelihood for preference data can be written as a function of policy log-ratios:
\[
\Pr\big(y^{(w)}\succ y^{(\ell)}\mid x\big)
=\sigmoid\!\Big(\beta\big[\log\tfrac{\pi^\star(y^{(w)}\mid x)}{\pi_{\mathrm{ref}}(y^{(w)}\mid x)}-\log\tfrac{\pi^\star(y^{(\ell)}\mid x)}{\pi_{\mathrm{ref}}(y^{(\ell)}\mid x)}\big]\Big).
\]
Replacing $\pi^\star$ by $\pi_\theta$ and maximizing the resulting log-likelihood yields exactly the DPO estimator:
\begin{equation}\label{eq:dpo}
\hat\theta_{\mathrm{DPO}}\in\arg\max_{\theta}\ \sum_{i=1}^N
\log \sigmoid\!\Big(\beta\big[\log\tfrac{\pi_\theta(y_i^{(w)}\mid x_i)}{\pi_{\mathrm{ref}}(y_i^{(w)}\mid x_i)}-\log\tfrac{\pi_\theta(y_i^{(\ell)}\mid x_i)}{\pi_{\mathrm{ref}}(y_i^{(\ell)}\mid x_i)}\big]\Big).
\end{equation}
\end{proposition}

\begin{proof}[Proof sketch, with the only nontrivial step highlighted]
Start from the BT model (A1):
\[
\Pr(y^{(w)}\succ y^{(\ell)}\mid x)=\sigmoid(r^\star(x,y^{(w)})-r^\star(x,y^{(\ell)})).
\]
Use (A2) to express reward differences by the log-ratio identity \eqref{eq:logratio}. The normalizer $\beta\log Z(x)$ cancels.
Substitute the resulting expression into the BT probability and take logs across observations. The optimization problem
in $\theta$ is precisely \eqref{eq:dpo}. The highlighted step is the inversion \eqref{eq:logratio}, which is a direct algebraic
consequence of the KL-regularized optimality condition.
\end{proof}

\noindent \textbf{Step 1: Invert the Policy to Express Rewards}

From the closed-form solution:
$$
\log \pi^*(y|x) = \frac{r^*(x,y)}{\beta} + \log \pi_{\text{ref}}(y|x) - \log Z(x)
$$

Rearrange to solve for $r^*(x,y)$:
$$
r^*(x,y) = \beta \log \frac{\pi^*(y|x)}{\pi_{\text{ref}}(y|x)} + \beta \log Z(x)
$$

\noindent \textbf{Step 2: Compute Reward Differences}

For any pair $(y_w, y_\ell)$:
\begin{align*}
r^*(x,y_w) - r^*(x,y_\ell) &= \beta \log \frac{\pi^*(y_w|x)}{\pi_{\text{ref}}(y_w|x)} + \beta \log Z(x) \\
&\quad - \left[\beta \log \frac{\pi^*(y_\ell|x)}{\pi_{\text{ref}}(y_\ell|x)} + \beta \log Z(x)\right] \\
&= \beta \left[\log \frac{\pi^*(y_w|x)}{\pi_{\text{ref}}(y_w|x)} - \log \frac{\pi^*(y_\ell|x)}{\pi_{\text{ref}}(y_\ell|x)}\right]
\end{align*}

\textbf{Observation:} The partition function $Z(x)$ cancels! This is \emph{crucial}—we don't need to compute the intractable normalization constant.

\noindent \textbf{Step 3: Substitute into the Bradley-Terry Likelihood}

The preference probability becomes:
\begin{align*}
P(y_w \succ y_\ell | x) &= \sigma(r^*(x,y_w) - r^*(x,y_\ell)) \\
&= \sigma\left(\beta \left[\log \frac{\pi^*(y_w|x)}{\pi_{\text{ref}}(y_w|x)} - \log \frac{\pi^*(y_\ell|x)}{\pi_{\text{ref}}(y_\ell|x)}\right]\right)
\end{align*}

\noindent \textbf{Step 4: Write the Parametric Log-Likelihood}

We don't observe $\pi^*$, so we posit a parametric family $\pi_\theta$ and maximize:

$$
\boxed{
\ell(\theta; \mathcal{D}) = \sum_{i=1}^N \log \sigma\left(\beta \left[\log \frac{\pi_\theta(y_i^{(w)}|x_i)}{\pi_{\text{ref}}(y_i^{(w)}|x_i)} - \log \frac{\pi_\theta(y_i^{(\ell)}|x_i)}{\pi_{\text{ref}}(y_i^{(\ell)}|x_i)}\right]\right)
}
$$

\textbf{This is exactly the DPO loss function.}

\noindent \textbf{Step 5: What Are We Estimating?}

\begin{itemize}
\item We are \emph{not} estimating $r^*$ directly
\item We are estimating the policy $\pi_\theta$ whose log-ratio scores best explain the preference data under the BT model
\item Under the structural assumption that this policy satisfies KL-regularized optimality, the estimated policy \emph{implicitly} encodes the reward function via the log-ratio identity
\end{itemize}

\begin{keypoint}
\textbf{The Isomorphism:} 
\begin{align*}
\text{MLE for } \pi_\theta \text{ under BT preferences} &\Leftrightarrow \text{Fitting the optimal policy under RLHF-KL}
\end{align*}

These are the same optimization problem. DPO \emph{is} structural estimation.
\end{keypoint}

\section{A one-page comparison: BT MLE vs.\ KL-regularized policy (and the isomorphism)}
The goal here is to let your brain see the mapping without doing any algebra.

\begin{center}
\renewcommand{\arraystretch}{1.25}
\begin{tabularx}{\textwidth}{@{}Y Y@{}}
\toprule
\textbf{Bradley--Terry (BT) MLE} & \textbf{KL-regularized tilted-Gibbs policy} \\\midrule
\textbf{Observable input} &
\textbf{Observable input} \\
Preference labels $(x, y^{(w)}, y^{(\ell)})$ & Reference policy $\pi_{\mathrm{ref}}(\cdot\mid x)$ and temperature $\beta$ \\
\midrule
\textbf{Latent object} &
\textbf{Latent object} \\
Score/utility $r^\star(x,y)$ entering choice probabilities & Reward $r^\star(x,y)$ entering planner's objective \\
\midrule
\textbf{Core assumption} &
\textbf{Core assumption} \\
$\Pr(y_1\succ y_2\mid x)=\sigmoid(r^\star(x,y_1)-r^\star(x,y_2))$ & $\pi^\star(\cdot\mid x)\in\arg\max_\pi \E_\pi[r^\star]-\beta\KL(\pi\|\pi_{\mathrm{ref}})$ \\
\midrule
\textbf{Optimization problem} &
\textbf{Optimization problem} \\
Maximize log-likelihood over parameters for $r_\eta$ or $\pi_\theta$ &
Maximize regularized expected reward over $\pi$ \\
\midrule
\textbf{Characterizing solution} &
\textbf{Characterizing solution} \\
FOC gives logistic regression in score differences (convex if linear) &
FOC gives softmax around reference: $\pi^\star \propto \pi_{\mathrm{ref}}\exp(r^\star/\beta)$ \\
\midrule
\textbf{Computing solution} &
\textbf{Computing solution} \\
Gradient ascent on $\sum\log\sigmoid(\Delta r)$ &
Closed form for $\pi^\star$ given $r^\star$; iterative methods only if $r^\star$ is implicit \\
\midrule
\textbf{Interpreting parameters} &
\textbf{Interpreting parameters} \\
Utilities identified up to additive constants (and scale) &
Policy trades reward vs.\ information/control cost; $\beta$ controls deviation from $\pi_{\mathrm{ref}}$ \\
\midrule
\textbf{Isomorphism (the bridge)} &
\textbf{Isomorphism (the bridge)} \\
Substitute $\Delta r^\star$ with policy log-ratio difference &
Invert the softmax: $\Delta r^\star=\beta\,\Delta\log\frac{\pi^\star}{\pi_{\mathrm{ref}}}$ \\
\midrule
\textbf{Resulting estimand} &
\textbf{Resulting estimand} \\
MLE over policy parameters $\theta$ using \emph{only} $\log(\pi_\theta/\pi_{\mathrm{ref}})$ &
Optimal policy for $r^\star$ can be expressed without $r^\star$ levels (only differences matter) \\
\bottomrule
\end{tabularx}
\end{center}

\section{So what is actually being ``estimated'' in DPO?}
There are three closely related but distinct answers. Confusing them causes most of the conceptual bugs.

\subsection{(1) Policy estimation: the direct object}
DPO is an $M$-estimator for the policy parameters $\theta$ in $\pi_\theta$ under a paired-comparison likelihood.
That is the literal statistical statement: you pick $\theta$ to maximize \eqref{eq:dpo}.

\subsection{(2) Utility differences: an implied object}
Given an estimated policy $\hat\pi$, you can \emph{define} an implied reward (up to constants)
\[
\hat r(x,y)\coloneqq \beta\log\frac{\hat\pi(y\mid x)}{\pi_{\mathrm{ref}}(y\mid x)} + \beta \hat c(x),
\]
where $\hat c(x)$ absorbs the unknown $\log Z(x)$. This recovers only differences in $r$, which is all BT identifies.

\subsection{(3) A structural control model: a maintained interpretation}
When you treat \eqref{eq:klproblem} as the \emph{structural} problem, DPO is estimating the policy that would be optimal
for some reward function under a KL control cost. The control model adds interpretation: $\pi_{\mathrm{ref}}$ is a prior,
$\beta$ is an information cost scale, and the log-ratio plays the role of a (generalized) ``value.'' But interpretation is only
as good as the maintained assumptions.

\section{Recursive methods and structural estimation: where the same algebra reappears}
In recursive macroeconomics (SLP, Ljungqvist-Sargent), we:
\begin{enumerate}
\item Define state space $\mathcal{S}$ and action space $\mathcal{A}$
\item Write the Bellman equation: $V(s) = \max_a [u(s,a) + \beta \mathbb{E}[V(s')]]$
\item Solve via value function iteration or policy iteration
\item Derive policy function: $\pi^*(s) = \arg\max_a [u(s,a) + \beta \mathbb{E}[V(s')]]$
\end{enumerate}

\textbf{The DPO framework extends this to stochastic policies with entropy regularization:}

\begin{theorem}[Soft Bellman Equation]
Define the soft value function:
$$
V^*(s) = \beta \log \sum_{a \in \mathcal{A}} \pi_{\text{ref}}(a|s) \exp\left(\frac{r(s,a) + \gamma \mathbb{E}[V^*(s')]]}{\beta}\right)
$$

The optimal policy is:
$$
\pi^*(a|s) = \frac{\pi_{\text{ref}}(a|s) \exp((r(s,a) + \gamma \mathbb{E}[V^*(s')])/\beta)}{\sum_{a'} \pi_{\text{ref}}(a'|s) \exp((r(s,a') + \gamma \mathbb{E}[V^*(s')])/\beta)}
$$
\end{theorem}

\textbf{Intuition:} The $\log \sum \exp$ is a ``soft max'' operator. As $\beta \to 0$, this converges to the standard Bellman equation. The reference policy $\pi_{\text{ref}}$ weights the alternatives.

\subsection{Why Gumbel Shocks Matter}

The Type-I Extreme Value (Gumbel) distribution is special because:

$$
\mathbb{E}\left[\max_a \{r(s,a) + \varepsilon_a\}\right] = \beta \log \sum_a \exp(r(s,a)/\beta) + \text{const}
$$

This is why:
\begin{itemize}
\item The soft Bellman equation has a log-sum-exp form
\item The choice probabilities are logit/softmax
\item Value function iteration remains a contraction
\end{itemize}

\begin{econintuition}
Gumbel shocks give us \emph{closed-form aggregation} of option values. The ``inclusive value'' or ``log-sum-exp'' term represents the expected value of the best choice under uncertainty. This same mathematics appears in:
\begin{itemize}
\item McFadden's logit demand
\item Rust's dynamic discrete choice
\item Soft Q-learning in RL
\item DPO for language model alignment
\end{itemize}
\end{econintuition}

\subsection{The Hotz-Miller Connection}

Hotz and Miller (1993) showed that in dynamic discrete choice models with logit shocks:

$$
V(s,a) - V(s,a') = \log P(a|s) - \log P(a'|s) + \log \nu(a) - \log \nu(a')
$$

where $\nu(a)$ are known functions of the shock distribution.

\textbf{DPO uses the same logic:}

$$
r^*(x,y) - r^*(x,y') = \beta \left[\log \frac{\pi^*(y|x)}{\pi_{\text{ref}}(y|x)} - \log \frac{\pi^*(y'|x)}{\pi_{\text{ref}}(y'|x)}\right]
$$

The reference policy $\pi_{\text{ref}}$ plays the role of the base measure, and $\beta$ is the scale parameter.


\section{Diagnostics and failure modes: when the equivalence fails (and how to notice)}
The equivalence is sharp. It breaks in predictable ways.

\begin{tcolorbox}[colback=light, colframe=accent, title=Checklist: what must be true for the algebra to go through]
\begin{itemize}
\item \textbf{(BT structure)} Preference labels must be representable as $\sigmoid(\Delta r)$ for some $r$.
\item \textbf{(KL control cost)} The planner must use KL divergence to the stated $\pi_{\mathrm{ref}}$ (not another regularizer).
\item \textbf{(Shared reward)} The same $r^\star$ must appear in both the preference model and the control problem.
\item \textbf{(Support)} $\pi_{\mathrm{ref}}>0$ on relevant outcomes so log-ratios exist.
\item \textbf{(Finite candidate implementation)} In practice, the estimand is defined on sampled candidates; interpret accordingly.
\end{itemize}
\end{tcolorbox}

\paragraph{Common subtlety: the candidate set matters.}
BT/logit is naturally a model over a fixed choice set. In LLM training, each $x$ effectively comes with a \emph{sampled} set of candidates
used for labeling. The cleanest interpretation is: you are estimating the policy restricted to that sampled set (or the induced distribution
over it), and extrapolation beyond it is model-dependent.

\paragraph{Rater heterogeneity.}
If different labelers have different reward functions $r^\star_j$, the BT model becomes a mixture. DPO then targets a compromise policy
that best fits the aggregated labels under the assumed single-$r^\star$ logistic model.

\section{Worked micro-example (finite candidate set)}
Fix a context $x$ and three candidates $\mathcal{Y}(x)=\{a,b,c\}$. Suppose a reference policy is
$\pi_{\mathrm{ref}}(\cdot\mid x)=(0.6, 0.3, 0.1)$ and the reward levels are $r^\star(x,\cdot)=(0,1,2)$ with $\beta=1$.
Then the optimal policy is
\[
\pi^\star(y\mid x)\propto \pi_{\mathrm{ref}}(y\mid x)e^{r^\star(x,y)}.
\]
Compute unnormalized weights:
$w_a=0.6e^0=0.6$, $w_b=0.3e^1\approx0.815$, $w_c=0.1e^2\approx0.739$,
so $\pi^\star\approx(0.28,0.38,0.34)$ after normalization.
The implied BT probability that $c$ beats $a$ is
$\sigmoid(r_c-r_a)=\sigmoid(2)\approx0.88$.
Using log-ratios instead:
\[
\beta\Big(\log\tfrac{\pi^\star(c)}{\pi_{\mathrm{ref}}(c)}-\log\tfrac{\pi^\star(a)}{\pi_{\mathrm{ref}}(a)}\Big)
= \log\tfrac{0.34}{0.1}-\log\tfrac{0.28}{0.6}\approx \log(3.4)-\log(0.47)\approx 1.22-(-0.76)=1.98,
\]
which matches $\Delta r=2$ up to rounding. Substitution reproduces the same BT probability. This is the equivalence in arithmetic form.

\newpage
\section{Application: Synthetic Cultural Agents}

\subsection{The Cultural Calibration Problem}

Standard LLMs are trained on WEIRD (Western, Educated, Industrialized, Rich, Democratic) data. To model non-WEIRD populations, we need to:

\begin{enumerate}
\item Define a cultural state vector $z$ parameterizing preferences
\item Construct a reward function $R(x,y;z)$ that varies with culture
\item Estimate a policy $\pi(y|x,z)$ that maximizes this reward
\end{enumerate}

\subsection{The GPS Framework}

Following Falk et al.\ (QJE 2018), we define:

$$
z = \begin{pmatrix}
\delta & \text{(time preference)} \\
\gamma & \text{(risk aversion)} \\
\xi^+ & \text{(positive reciprocity)} \\
\xi^- & \text{(negative reciprocity)} \\
\alpha & \text{(altruism)} \\
\tau & \text{(trust)}
\end{pmatrix}
$$

These are not ad-hoc labels but \emph{structural parameters} estimated from incentivized experiments in 76 countries.

\subsection{The Synthetic Ethnography Pipeline}

\textbf{Problem:} For many populations, we lack digital trace data to train models directly.
\\
\textbf{Solution:} Use the DPO isomorphism to create synthetic preference data:

\begin{enumerate}
\item \textbf{Teacher Model:} Condition a frontier instruction model on ethnographic profiles and GPS parameters
\item \textbf{Generate Synthetic Preferences:} Produce $(x, y_w, y_\ell)$ triplets reflecting cultural preferences
\item \textbf{Student Model:} Use DPO to distill these preferences into a smaller, open-source model (e.g., nemotron, mixtral, gpt-oss)
\item \textbf{Validation:} Check that the student's behavior matches held-out moments from WVS, GPS, etc.
\end{enumerate}

\subsection{Why This Works}

The DPO isomorphism means that:
\begin{itemize}
\item Synthetic preference data \emph{encodes} the reward function $R(x,y;z)$
\item The student learns the optimal policy $\pi^*(\cdot|x,z)$ without needing an explicit reward model
\item Cultural parameters $z$ shift the induced policy in a structurally interpretable way
\end{itemize}

This is \emph{transfer learning as structural estimation}—we're using a frontier model's world knowledge to calibrate behavioral parameters for populations outside its training distribution.

\newpage
\section{Conclusion and Further Directions}

\subsection{What We've Learned}

\begin{enumerate}
\item \textbf{Optimal policies and MLE estimators are the same object} under revealed preference logic with KL-regularized optimality.

\item \textbf{DPO is a structural method}, not an ad-hoc ML trick. It recovers policies via inversion of choice probabilities—exactly the strategy used in structural DDC.

\item \textbf{The Bradley-Terry model and KL-regularized RL are isomorphic} when rewards are expressed as log-ratios relative to a reference policy.

\item \textbf{Recursive methods and structural estimation connect} through the Hotz-Miller inversion: both avoid repeatedly solving fixed-point problems by exploiting closed-form relationships between policies and values.

\item \textbf{Cultural alignment can be formalized} using GPS-calibrated reward functions and synthetic ethnography via teacher-student distillation.
\end{enumerate}

\subsection{Open Questions for Further Study}

\begin{enumerate}
\item \textbf{Misspecification:} What happens when preferences violate BT (e.g., intransitivity, context effects)? DPO remains a valid M-estimator, but what is the \emph{economic interpretation} of the fitted policy?

\item \textbf{Dynamics:} The one-step analysis extends to sequential decision-making via soft Bellman equations. How do we estimate \emph{recursive} cultural preferences over trajectories?

\item \textbf{Heterogeneity:} Real populations have distributions over $z$. Can we extend DPO to mixture models or hierarchical specifications?

\item \textbf{Counterfactuals:} Once we have $\hat{\pi}(y|x,z)$, how do we conduct policy analysis? What is the ``treatment effect'' of changing cultural parameters?

\item \textbf{Identification:} What variation in the data identifies different components of $z$? This is a classic structural IV problem.
\end{enumerate}
\subsection{Suggested Readings}

\textbf{Structural Econometrics:}
\begin{itemize}
\item Hotz \& Miller (1993): ``Conditional Choice Probabilities and the Estimation of Dynamic Models''
\item Rust (1987): ``Optimal Replacement of GMC Bus Engines''
\item Train (2009): \emph{Discrete Choice Methods with Simulation}
\end{itemize}

\noindent \textbf{Recursive Methods:}
\begin{itemize}
\item Stokey, Lucas \& Prescott (1989): \emph{Recursive Methods in Economic Dynamics}
\item Ljungqvist \& Sargent (2018): \emph{Recursive Macroeconomic Theory}
\end{itemize}

\noindent \textbf{Machine Learning \& Alignment:}
\begin{itemize}
\item Rafailov et al.\ (2023): ``Direct Preference Optimization: Your Language Model is Secretly a Reward Model''
\item Falk et al.\ (2018): ``Global Evidence on Economic Preferences''
\item Ouyang et al.\ (2022): ``Training Language Models to Follow Instructions with Human Feedback''
\end{itemize}

\vfill
\end{document}
